\documentclass[11pt,a4paper]{article}

\usepackage{amsmath}
\usepackage{amssymb}
\usepackage{amsthm}
\usepackage{booktabs}
\usepackage{array}
\usepackage{geometry}
\usepackage{enumitem}
\usepackage[hidelinks]{hyperref}

\geometry{margin=1in}

\theoremstyle{plain}
\newtheorem{theorem}{Theorem}[section]
\newtheorem{lemma}[theorem]{Lemma}
\newtheorem{proposition}[theorem]{Proposition}
\newtheorem{corollary}[theorem]{Corollary}
\theoremstyle{definition}
\newtheorem{definition}[theorem]{Definition}
\newtheorem{hypothesis}[theorem]{Hypothesis}
\theoremstyle{remark}
\newtheorem{remark}[theorem]{Remark}

\newcommand{\GD}{G^{D}}
\newcommand{\ED}{E^{D}}
\newcommand{\orient}{\operatorname{or}}
\newcommand{\Cone}{\operatorname{Cone}}

\title{Two forbidden patterns bound the boundary intersections\\
of two simple heptagons\\[2mm]
\large A computer-assisted proof that $f(7,7)=38$,\\
modulo one imported general-position reduction}

\author{Jayden Lin\thanks{Northeastern University. This work was carried out
with substantial generative-AI assistance; see \S\ref{sec:ai}.}}

\date{17 September 2026}

\begin{document}
\maketitle

\begin{abstract}
Let $f(m,n)$ be the maximum number of points at which the boundaries of two
simple polygons with $m$ and $n$ sides can meet. Dillencourt, Mount and
Saalfeld posed the problem in 1993 and settled the cases where $m$ or $n$ is
even; the odd--odd cases have resisted. We give a computer-assisted proof that
$f(7,7)=38$, conditional on one imported statement.

The mathematical content is a reduction. Rather than enumerate the
$13{,}003{,}620$ candidate crossing patterns with $40$ intersections, we
exhibit \emph{two} local patterns, involving $8$ and $14$ prescribed
side-pair incidences, and show that every candidate pattern with $40$
intersections and every one with $42$ contains one of them. Each of the two is
non-realizable, and a degree argument over $28$ cases reduces the whole
exclusion to those two facts.

Everything in this paper except one statement is verified by the Lean~4 kernel
from primitive definitions, with no appeal to any published lemma. The
exception is the reduction to general position, for which we import
Ackerman--Keszegh--Rote \S9 and carry it as an explicit hypothesis. We are
explicit throughout about which assertions are machine-checked, which are
imported, and which --- notably the faithfulness of our definitions --- remain
matters for a human reader.
\end{abstract}

\tableofcontents

\section{Introduction}

\subsection{The problem}

Two simple polygons in the plane, with $m$ and $n$ sides, are in
\emph{general position} if no vertex of either lies on the boundary of the
other and no three of the combined $m+n$ vertices are collinear. Write
$f(m,n)$ for the maximum, over such pairs, of
$\lvert \partial P \cap \partial Q\rvert$.

\begin{quote}
\textbf{Problem} (Dillencourt--Mount--Saalfeld, 1993). Determine $f(m,n)$.
\end{quote}

They determined $f(m,n)$ when at least one of $m,n$ is even. The odd--odd cases
are open. A 2012 preprint of G\"unther claiming $f(n,n)$ for odd $n$ was
withdrawn by its author, with the stated reason that a counterexample was found
to a weakened form of one of its lemmas. The smallest open case is
$m=n=7$.

Two elementary facts frame it. Each side of $P$ meets each side of $Q$ at most
once in general position, so $I\le 49$; and since $\partial P$ and $\partial Q$
are closed curves meeting transversally, $I$ is even. The published state of the
art gives $I\le 40$, and $38$ is attained. So the entire question is whether
$40$ is attainable.

\subsection{What is proved here}

\begin{theorem}[Main, conditional]\label{thm:main}
Let $P$ and $Q$ be simple heptagons in the plane, each simple as a point set,
such that no vertex of either lies on the boundary of the other. Then
\[
  \lvert \partial P \cap \partial Q\rvert \;\le\; 38 ,
\]
and $38$ is attained. No condition on collinearity of vertices is assumed.
\end{theorem}

The hypothesis ``no vertex of either lies on the boundary of the other'' is
exactly the conclusion of Ackerman--Keszegh--Rote \S9. We do not prove it;
\S\ref{sec:import} states precisely what we import and why, and
\S\ref{sec:e9attempt} reports an unsuccessful attempt to remove it.

\begin{remark}
Theorem~\ref{thm:main} is \textbf{not} $f(7,7)=38$, and we are careful not to
write it that way anywhere. It is the statement one obtains after the
general-position reduction has been applied, and the reduction is imported.
\end{remark}

\subsection{What is new}

The contribution is a reduction, not a technique. Both ingredients we use ---
encoding realizability of a crossing pattern as an oriented-matroid
satisfiability problem, and refuting it mechanically --- are standard; see
\S\ref{sec:related}. What appears to be new is:

\begin{itemize}[leftmargin=*]
  \item \textbf{Two patterns suffice} (Theorem~\ref{thm:cover}). A search that
    was $13{,}003{,}620$ cases becomes two non-realizability facts plus a
    degree argument over $28$ cases.
  \item \textbf{The smaller pattern is human-readable} and needs no simplicity
    hypothesis: it constrains three consecutive sides of each polygon and
    prescribes eight incidences (\S\ref{sec:patterns}).
  \item \textbf{The reduction to the plane is unconditional.} The combinatorial
    argument is about orientation signs; getting from there to genuine point
    sets in $\mathbb{R}^2$ --- including removing the collinearity assumption
    that \textsc{akr} \S3 asserts without proof --- is carried out in full
    (\S\ref{sec:plane}).
\end{itemize}

A bounded prior-art search found no previous statement of the covering result.
We stress in \S\ref{sec:related} that this is weak evidence and not a novelty
claim.

\section{Preliminaries}\label{sec:prelim}

\subsection{Configurations and the disjointness graph}

Let $P$ have vertices $p_0,\dots,p_6$ and $Q$ have $q_0,\dots,q_6$, indices mod
$7$, with sides $P_i=[p_i,p_{i+1}]$ and $Q_j=[q_j,q_{j+1}]$. For points
$a,b,c$ write
\[
  \orient(a,b,c)=(b_1-a_1)(c_2-a_2)-(b_2-a_2)(c_1-a_1),
\]
twice the signed area. Sides $[a,b]$ and $[c,d]$ \emph{cross} when
\[
  \orient(a,b,c)\cdot\orient(a,b,d)<0
  \quad\text{and}\quad
  \orient(c,d,a)\cdot\orient(c,d,b)<0 .
\]
This is a conjunction of \emph{strict} inequalities, a fact used repeatedly: it
makes crossing an open condition.

\begin{definition}
The \emph{disjointness graph} $\GD$ is the bipartite graph on
$\{P_0,\dots,P_6\}\sqcup\{Q_0,\dots,Q_6\}$ with an edge $P_iQ_j$ when $P_i$ and
$Q_j$ do \emph{not} cross.
\end{definition}

In general position each side pair contributes at most one intersection point,
and exactly one when it crosses, so writing $I$ for the number of intersection
points and $\lvert \ED\rvert$ for the number of edges of $\GD$,
\begin{equation}\label{eq:link}
  I + \lvert \ED\rvert = 49 .
\end{equation}

\subsection{Two unconditional bounds}

\begin{lemma}\label{lem:even}
$I$ is even.
\end{lemma}

\begin{lemma}\label{lem:42}
$\delta(\GD)\ge 1$, and consequently $I\le 42$.
\end{lemma}

\begin{proof}[Proof sketch]
Every side of $P$ fails to cross at least one side of $Q$ and vice versa, so
$\lvert \ED\rvert\ge 7$; with \eqref{eq:link} this gives $I\le 42$.
\end{proof}

\begin{remark}
\v{C}ern\'y, Kyn\v{c}l and T\'oth obtain $\delta(\GD)\ge1$ from the fact that a
line meets a simple $7$-gon in at most $6$ sides. We do not import that: the
minimum-degree statement is derived directly and is machine-checked
unconditionally, so their Lemma~1 is corroboration here, not a dependency.
\end{remark}

By Lemmas~\ref{lem:even} and~\ref{lem:42}, $I\in\{\dots,38,40,42\}$, and since
$38$ is attained (\S\ref{sec:lower}) the theorem reduces to excluding $40$ and
$42$.

\subsection{Orientation data}

All of \S\S\ref{sec:patterns}--\ref{sec:exclude} concerns only the signs of the
$\binom{14}{3}=364$ determinants $\orient(v_a,v_b,v_c)$ on the combined vertex
set. A \emph{configuration} is a $14$-tuple of points together with the
requirement that all $364$ signs be nonzero; a \emph{crossing pattern} is the
resulting $7\times7$ Boolean matrix of side-pair crossings. Realizability
questions about patterns are therefore questions about which sign vectors arise
from actual point sets --- that is, about chirotopes.

\section{Two forbidden patterns}\label{sec:patterns}

\subsection{The small pattern}

\begin{definition}\label{def:pat8}
Pattern $\Pi_8$ constrains sides $P_0,P_1,P_2$ and $Q_0,Q_1,Q_2$:
\[
\begin{array}{c|ccc}
      & Q_0 & Q_1 & Q_2 \\ \hline
  P_0 & \times & \times & \times \\
  P_1 & \times & \cdot   & \times \\
  P_2 & \times & \times & \ast
\end{array}
\]
where $\times$ means \emph{crosses}, $\cdot$ means \emph{does not cross}, and
$\ast$ is unconstrained. That is eight prescribed incidences: seven crossings
and one non-crossing.
\end{definition}

\begin{theorem}\label{thm:pat8}
No configuration realizes $\Pi_8$.
\end{theorem}

Theorem~\ref{thm:pat8} assumes nothing beyond the eight incidences and the
nondegeneracy of the $364$ determinants: in particular it needs \emph{no}
hypothesis that $P$ or $Q$ is simple. It is the single most compact statement
in this work, and we suggest it as the first thing a sceptical reader should
check.

\subsection{The larger pattern}

\begin{definition}\label{def:pat14}
Pattern $\Pi_{14}$ constrains $P_0,\dots,P_3$ and $Q_0,\dots,Q_4$ by fourteen
incidences,
\[
\begin{array}{c|ccccc}
      & Q_0 & Q_1 & Q_2 & Q_3 & Q_4 \\ \hline
  P_0 & \ast & \times & \times & \ast & \times \\
  P_1 & \times & \cdot & \times & \times & \times \\
  P_2 & \times & \times & \ast & \ast & \times \\
  P_3 & \times & \times & \times & \ast & \ast
\end{array}
\]
together with seven simplicity conditions: $P_0\!\not\times\!P_2$,
$P_0\!\not\times\!P_3$, $P_1\!\not\times\!P_3$, $Q_0\!\not\times\!Q_2$,
$Q_0\!\not\times\!Q_3$, $Q_1\!\not\times\!Q_3$, $Q_1\!\not\times\!Q_4$.
\end{definition}

\begin{theorem}\label{thm:pat14}
No configuration realizes $\Pi_{14}$.
\end{theorem}

Unlike $\Pi_8$, this one genuinely uses that the polygons do not cross
themselves; the seven self-non-crossing conditions are part of the hypothesis
and are supplied by simplicity.

\subsection{How the patterns are proved non-realizable}

Both theorems are established by encoding realizability as a Boolean
satisfiability problem over the $364$ orientation variables --- Grassmann--Pl\"ucker
relations, acyclicity, and the prescribed incidences --- and refuting it. The
refutation is not trusted:

\begin{enumerate}[leftmargin=*]
  \item Each clause of the encoding is \emph{derived from the geometry} inside
    the proof assistant. The Grassmann--Pl\"ucker and acyclicity clauses are
    polynomial identities in the coordinates; the pattern clauses are the
    definition of crossing unfolded. Nothing about the encoding is assumed.
  \item The solver's refutation is checked, then transcribed into an explicit
    chain of unit-propagation steps. The kernel verifies a chain of
    implications between propositions; \emph{no search runs inside the
    proof assistant}.
\end{enumerate}

\begin{remark}[Positive control]
A refutation that also covered a configuration which exists would be proving
something false. Both patterns require $P_1$ to miss $Q_1$; in the
$38$-crossing configuration of \S\ref{sec:lower} those two sides cross, so
neither pattern applies to it. This control is itself machine-checked.
\end{remark}

\section{The covering theorem}\label{sec:cover}

\begin{theorem}\label{thm:cover}
Every crossing pattern of two simple heptagons with $I=40$, and every one with
$I=42$, contains $\Pi_8$ or $\Pi_{14}$ after relabelling.
\end{theorem}

Here ``after relabelling'' refers to the group of order $392$ generated by
rotating and reflecting each polygon and exchanging $P$ with $Q$.

\subsection{The $I=42$ case}

By \eqref{eq:link}, $I=42$ forces $\lvert \ED\rvert=7$, and with
$\delta(\GD)\ge1$ on both sides the graph is a perfect matching. There are
$7!=5040$ of these, all equivalent under relabelling to few orbits, and each
contains one of the two patterns.

\subsection{The $I=40$ case}

Here $\lvert \ED\rvert=9$, so $\GD$ has nine edges with all fourteen degrees at
least~$1$. Write $r_i=\deg P_i$ and $c_j=\deg Q_j$, so
$\sum_i r_i=\sum_j c_j=9$ with every $r_i,c_j\ge1$. Since nine edges are spread
over seven rows each of degree at least one, the multiset $\{r_i\}$ is one of
few shapes, and likewise for columns; enumerating gives \textbf{28 degree
shapes}. For each shape a short argument locates a sub-pattern equivalent to
$\Pi_8$ or $\Pi_{14}$.

The point of organising the argument this way is that the $28$ cases are a
human-sized object, whereas the $13{,}003{,}620$ patterns are not.

\section{Excluding 40 and 42}\label{sec:exclude}

\begin{corollary}\label{cor:not4042}
No configuration has $I=40$ or $I=42$.
\end{corollary}

\begin{proof}
Immediate from Theorems~\ref{thm:pat8}, \ref{thm:pat14} and~\ref{thm:cover}.
\end{proof}

\begin{corollary}\label{cor:38}
Every configuration has $I\le 38$.
\end{corollary}

\begin{proof}
$I\le42$ by Lemma~\ref{lem:42}, $I$ is even by Lemma~\ref{lem:even}, and $40$
and $42$ are excluded by Corollary~\ref{cor:not4042}.
\end{proof}

\section{The lower bound}\label{sec:lower}

\begin{proposition}\label{prop:lower}
The heptagons with vertices
\[
\begin{aligned}
P &: (373,157),\,(132,5),\,(369,182),\,(88,157),\,(327,282),\,(60,272),\,(372,298),\\
Q &: (344,24),\,(89,303),\,(345,19),\,(373,214),\,(270,338),\,(325,93),\,(192,284)
\end{aligned}
\]
are simple, are in general position, and their boundaries meet in exactly $38$
points.
\end{proposition}

All three assertions are decided by exact integer arithmetic and verified by
the kernel. Combined with Corollary~\ref{cor:38} this gives the equality in
Theorem~\ref{thm:main}.

\section{From orientation signs to the plane}\label{sec:plane}

Sections~\ref{sec:patterns}--\ref{sec:exclude} are about sign vectors. Three
further steps are needed to reach Theorem~\ref{thm:main}, and all three are
carried out unconditionally.

\subsection{Integer coordinates are no restriction}

The combinatorial development uses integer coordinates. This costs nothing:
$\orient$ is homogeneous of degree~$2$, so the set of point-tuples with a given
sign vector is invariant under scaling. Scaling a configuration until a box
around it has side length exceeding~$1$ produces an integer point in that box
with the same $364$ signs. No rationals and no clearing of denominators are
required.

\subsection{Counting points, not side pairs}

The combinatorial argument counts crossing \emph{pairs}; the theorem counts
\emph{points}. In general position these agree: a crossing pair meets in
exactly one point, and a point determines its pair. Both directions are proved.

\subsection{Removing the collinearity assumption}\label{sec:removecollin}

\textsc{akr} \S3 assumes, ``for convenience'' and without proof, that no three
of the combined vertices are collinear. We do not assume it. The argument is a
perturbation, in two halves:

\begin{lemma}\label{lem:open}
Crossing is an open condition, so for a small enough perturbation of the
fourteen vertices no existing crossing is lost.
\end{lemma}

\begin{lemma}\label{lem:simple}
If two segments are disjoint, they remain disjoint when their four endpoints
move by less than some $\varepsilon>0$.
\end{lemma}

Lemma~\ref{lem:open} handles the crossings; Lemma~\ref{lem:simple} handles
simplicity, which is a \emph{closed} condition and therefore out of reach of any
openness argument. Since the configurations in general position are dense,
combining the two lemmas approximates an arbitrary pair of point-set simple
heptagons by one in general position with at least as many crossings, giving
the bound for the original. Lemma~\ref{lem:simple} is proved for an arbitrary
real normed space and is independent of everything else here.

\begin{remark}
Note the direction: we need the perturbed configuration to have \emph{at least}
as many crossings, not the same number, and openness gives exactly that.
\end{remark}

\section{What is imported}\label{sec:import}

\begin{hypothesis}[\textsc{(E9)}]\label{hyp:e9}
No vertex of either polygon lies on the boundary of the other.
\end{hypothesis}

This is the conclusion of \textsc{akr} \S9, where \S9.1 removes vertex--vertex
coincidences and \S9.2 vertex--side ones. It is the \emph{only} external
mathematical input to Theorem~\ref{thm:main}.

Why it is needed, precisely: without it, a point of
$\partial P\cap\partial Q$ can lie on \emph{no} crossing pair. Two sides that
merely touch share a point but have vanishing determinants, so the pair is not
counted, and the bound on pairs does not bound points. The failure is in the
unfavourable direction, which is why \S\ref{sec:removecollin} does not also
dispose of it.

\begin{remark}[Citation]
It would be an under-citation to attribute all of the usual general-position
package to \S9. \S9 gives Hypothesis~\ref{hyp:e9} and nothing more; the
collinearity clause has its only warrant in one unproved sentence of \S3, and
is discharged here (\S\ref{sec:removecollin}). Additionally, that the perturbed
polygons remain simple is stated nowhere in that paper. We fold simplicity into
Hypothesis~\ref{hyp:e9} deliberately, so that the gap is visible.
\end{remark}

\subsection{An attempt to remove it}\label{sec:e9attempt}

We attempted to prove Hypothesis~\ref{hyp:e9} and did not succeed. The
decomposition into its vertex--vertex and vertex--side halves is formalized, as
are two local lemmas, but the global reduction is not. Two negative
observations came out of the attempt and are worth recording, because they
close off the obvious routes:

\begin{itemize}[leftmargin=*]
  \item Generic perturbation preserving proper crossings does \emph{not} imply
    that the point count is monotone.
  \item Naive displacement into the interior of a cone does not either.
\end{itemize}

Explicit coordinates exhibiting both failures are recorded in the project
repository; they are analytic and \emph{not} machine-checked. This is
consistent with the known bad case $\Cone_P(v)\subseteq\Cone_Q(w)$, on which an
earlier attempt of ours broke. None of this is evidence that \textsc{akr}'s
theorem is false, and we allege no error in it; what remains is substantial
unfinished work.

\section{Machine verification, and its limits}\label{sec:machine}

Every mathematical assertion in \S\S\ref{sec:prelim}--\ref{sec:plane} is
verified by the Lean~4 kernel. There is no \texttt{sorry}, no
\texttt{native\_decide}, and no project-specific axiom; every theorem depends on
at most the three standard axioms \texttt{propext},
\texttt{Classical.choice}, \texttt{Quot.sound}. Hypothesis~\ref{hyp:e9} is a
hypothesis of the statement, never an axiom.

\begin{remark}[What the kernel does not check]\label{rem:faithful}
The kernel proves theorems about definitions; it does not certify that the
definitions mean what we say. That $\partial P\cap\partial Q$ is modelled by the
right set, and that ``simple heptagon'' is modelled by the right predicate,
are judgements a human reader must make. This is, in our view, the principal
remaining risk, and it is not reducible by more computation.

One partial check is available and we have carried it out: the definitions,
applied to the explicit $38$-crossing configuration of
Proposition~\ref{prop:lower}, yield exactly $38$. A wrong definition would have
to be wrong in a way that still lands on that value.
\end{remark}

\subsection{Independent corroboration}

Beyond the kernel: the terminal case analysis of an earlier version of this
argument was reproduced by a second, independently written implementation in a
different language; all solver refutations are re-checked by a standalone
proof checker, which is itself first shown to reject non-refutations; and all
generated artifacts are regenerated and compared byte-for-byte. These guard
against implementation error. \emph{None of them is a check on whether the
argument is the right argument.}

\section{Related work}\label{sec:related}

Dillencourt, Mount and Saalfeld posed the problem and settled the even cases.
\v{C}ern\'y, Kyn\v{c}l and T\'oth gave bounds for odd cases and solved $k=5$,
explicitly leaving $k=\ell=7$ open. Ackerman, Keszegh and Rote obtained
asymptotically optimal bounds; their \S9 is what we import. G\"unther's 2012
preprint on odd $n$ was withdrawn by its author.

Both of our tools are standard. Encoding order-type or chirotope realizability
as a satisfiability problem and refuting it mechanically is well established,
including in combination with proof assistants. Framing extremal results in
terms of forbidden configurations is likewise standard. What we have not found
in the literature is either applied to crossing patterns of two simple
polygons, or to $f(m,n)$.

\begin{remark}[On novelty]
We searched for prior statements of the covering result and found none. We
record plainly that this is \textbf{weak evidence and not a novelty claim}: the
search was bounded, the journal version of one central reference was not
reachable, the forward citations of another were not examined, and the relevant
literature is small enough that a result could sit in it under different
vocabulary. A domain expert's judgement is not replaced by such a search.
\end{remark}

\section{Status, and what we ask of a reader}\label{sec:status}

To state the position exactly:

\begin{itemize}[leftmargin=*]
  \item $f(7,7)=38$ is \textbf{not} proved. Theorem~\ref{thm:main} carries
    Hypothesis~\ref{hyp:e9}.
  \item Everything else is machine-checked from primitive definitions with no
    published lemma imported.
  \item \textbf{No mathematician outside this process has reviewed any of it.}
\end{itemize}

The third point is the binding constraint, and it is why this paper exists. Two
entry points need the least context:

\begin{enumerate}[leftmargin=*]
  \item \textbf{Theorem~\ref{thm:pat8}.} Eight incidences on three consecutive
    sides of each polygon, no simplicity hypothesis. If it is wrong, everything
    here is wrong.
  \item \textbf{Lemma~\ref{lem:simple}.} A self-contained statement about
    segments in a normed space, whose formalization depends on nothing else in
    the development.
\end{enumerate}

We would also welcome an opinion on Remark~\ref{rem:faithful} --- whether our
formal definitions capture the intended objects --- and on whether the covering
result is in fact already known.

\section*{Disclosure of generative-AI use}\label{sec:ai}
\addcontentsline{toc}{section}{Disclosure of generative-AI use}

Generative-AI systems were used substantially throughout: exploring the
reduction, generating essentially all software including the formalization,
auditing the cited sources, and drafting this paper. The named author is
responsible for the claims. AI output was not treated as trustworthy: each
major claim is discharged by a separately recorded mechanism --- source audit,
independent reimplementation, a machine-checkable certificate with a standalone
checker, or kernel verification --- and errors introduced during development
are catalogued in the project repository rather than quietly repaired. The
patterns of \S\ref{sec:patterns} were \emph{found} by computer search; what they
are is small enough to check by hand.

\end{document}
